\section{Blank initialization, unbounded runtime, and the ZFC bridge}
\label{sec:runtime235}

The literal fold theorem connects the new machine to its rotated binary
basis. The source graph theorem connects the new register program to the
inherited register program. A third theorem is needed between the two:
the exact binary tables must execute the transformed source on every
reachable runtime frame. This section describes that theorem and the
finite initialization which supplies its first frame.

\subsection{The exact initial runtime frame}

Starting on a blank tape at head coordinate 0, all three checked machines
reach inherited dispatcher state 96 after exactly 409,835 transitions.
The head is at $-5$, the marker is at $-6$, and the register-file header
is at coordinate 176. Exactly 351 tape cells contain a 1; all other cells
are zero. The entire sorted support is in \texttt{evidence/startup235.json}.
The compact JSON serialization of that support, including a final newline,
has SHA-256
\begin{center}\small\ttfamily
2f1158059a8021e206c790eeb43509d0b2c4bb3b73471ab47edcb73242e61772.
\end{center}

The logical register tuple is
\[
 R_{\mathrm{seed}}=(3,0,0,0,3,0,0,0,0).
\]
There are 167 spare blocks, each with value 1, and the extra-gap parameter
is 117. The logical-to-physical register map is
\[
 (0,1,2,3,4,5,6,7,8)\longmapsto(6,2,0,4,3,7,8,5,1).
\]
Consequently the primary physical values are
$(0,0,0,3,0,0,3,0,0)$. The pinned ancestor's seed check and the new
whole-tape equality identify precisely the calibrated source seed;
matching only the number of one-cells would not suffice.

The independent startup checker uses an unbounded sparse tape for the
pinned reference and two dense tapes for the rotated basis and an
independently canonicalized 235-state machine. At each of the 409,836
checkpoints it verifies the state correspondence, head, read symbol and
literal transition. Since the complete tapes begin equal and every
write agrees, this is an inductive equality of whole configurations.
It also compares the complete final supports directly. The head stays
within $[-15,701]$, strictly inside both dense allocations, so the finite
arrays do not approximate an unvisited boundary. No changed source row
and neither of the two folded states is visited in initialization.

The seed has an empty source path. This needs no extra initial decoder
lemma: the first leaf of a complete binary tree is an all-zero word.
That word with its padding is literally the same tape word as the empty
path with 64 padding zeros. The initialized configuration is already an
instance of the first guarded runtime case.

\subsection{The quantified instruction interface}

Write $B(v)=0\,1^{v+1}$ for a unary register block. Relative to the
dispatcher head, a frame consists of a marker 1 immediately to its left,
then
\[
 p\,0^{64+g-|p|}\,1\;
 B(a_0)\cdots B(a_8)\;B(s_0)B(s_1)\cdots B(s_k),
\]
with zero exterior in both directions. Here $p$ is a source address,
$g\ge0$, the nine $a_i$ are the physical primary values, and $k\ge0$.
Thus at least one spare block exists. All register values and spare
values are arbitrary natural numbers. Zero padding makes a prefix
address and its leftmost descendant address denote the same frame.

For each increment instruction the required endpoint increments its
selected logical register, appends a spare block of value 1, and advances
the source address. For a positive decrement it removes one unit and
advances; for a zero decrement it follows the specified break edge.
Break and continue instructions preserve registers and select their
structured successor. A halt instruction reaches outcome 254. The
endpoint includes the exact marker, padded address, gap adjustment,
entire primary file, complete spare list, head displacement and infinite
zero exterior for every nonhalt instruction. A halt instruction is checked
for arrival at outcome 254 with positive duration; no final tape or head
condition is needed at that terminal outcome. These are stronger contracts
than register-value tests.

There are 6,455 source leaves. Each decrement has separate zero and
positive cases, yielding exactly 8,386 guarded cases. In a positive
case, the selected value is represented by $1+n$ with arbitrary
$n\in\mathbb N$; in a zero case it is exactly zero. All other primary
values remain independent arbitrary naturals. The initialization and
every successful instruction preserve the frame class, so induction
composes these local contracts over every finite source run.

\subsection{Why the runtime certificates are unbounded}

The retained engine has only three kinds of executable inference.
Ordinary transitions perform a literal row of the exact selected table.
A unary self-loop in the current scan direction is accelerated by
induction on an arbitrary natural run length. A list traversal is
accelerated only after the corresponding one-block theorem has been
rederived on this exact table in arbitrary exterior contexts.

Put $C(v)=1\,0\,1^v$. The four required one-block schemas are:
\begin{center}
\begin{tabular}{rccl}
\toprule
Entry state & Direction & Transformation & Exact duration\\
\midrule
73 & $R$ & $B(v)\to B(v)$ & $v+2$\\
95 & $R$ & $B(v)\to C(v)$ & $v+2$\\
48 & $L$ & $B(v)\to B(v)$ & $v+2$\\
112 & $L$ & $C(v)\to B(v)$ & $v+2$\\
\bottomrule
\end{tabular}
\end{center}
The list cases follow by induction on a finite list of such blocks,
including the empty list. The normalization identity
$C(A)\,1=1\,B(A)$ for a block list $A$ is itself the concatenation of
the one-block identity $C(v)1=1B(v)$; the empty case and its reversal
are included. No finite sample of list lengths is substituted for this
induction.

The symbolic tape represents unary lengths by affine expressions in
13 slots: a constant, the extra gap, nine logical primary values, the
first spare value, and the total width of the remaining spare list.
It retains the opaque list structure where structural induction is
required. Every successful instruction produces a nonnegative affine
duration with strictly positive constant term. Thus, for every concrete
assignment of natural values and every finite spare list, its macro is
finite and takes at least one actual transition.

The engine permits at most 4,000 symbolic rule applications while
\emph{constructing one proof}. Reaching that limit produces a failure,
not an assumed theorem. A successful rule can represent arbitrarily many
literal transitions, so this budget bounds neither the actual run time
nor the register values, gap, or spare-list length. A symbolic expression
that could be empty cannot be read as a literal symbol without an
appropriate rule; this check rejected several false candidates.

The new driver pins the input source, literal table, retained induction
engine and helper by SHA-256. It independently rechecks the complete
source graph, derives the four list schemas anew, and proves every guarded
case. The final results are:
\begin{center}
\begin{tabular}{lrrrr}
\toprule
Literal table & Cases & Symbolic rules & Rows used & Min. duration\\
\midrule
Rotated 236-state basis &8,386&2,266,171&472/472&13\\
Final 235-state machine &8,386&2,269,055&470/470&13\\
\bottomrule
\end{tabular}
\end{center}
Every literal row is used by at least one proved case. Row coverage is a
useful census, but the proof is the exact tape and successor identity for
each nonhalt symbolic case, and the halt endpoint for a halt case.
Both routes are delivered: the basis proof plus the
independent fold theorem, and a direct fresh proof of the final table.
The case arrays record durations, symbolic lengths and rule-kind counts.
Their full generation is reproducible from the retained literal engine.

\subsection{Composing the logical equivalence in ZFC}

Let $S_{236}$ be the pinned register source, $S_7$ the seven-rotation
source, $B_{236}$ the new binary basis, and $M_{235}^{\mathrm{old}}$ the
folded machine with inherited labels. The established equivalences are
\begin{align*}
 \Halt(M_{235},\mathrm{blank})
 &\iff \Halt(M_{235}^{\mathrm{old}},\mathrm{blank})
       &&\text{exact state renumbering},\\
 &\iff \Halt(B_{236},\mathrm{blank})
       &&\text{positive 1/3-step fold},\\
 &\iff \Halt(S_7,R_{\mathrm{seed}})
       &&\text{startup and runtime induction},\\
 &\iff \Halt(S_{236},R_{\mathrm{seed}})
       &&\text{complete source graph equality},\\
 &\iff \Refut(\ZFC)
       &&\text{inherited logical calibration}.
\end{align*}
The last line retains the existing source-to-native-refutation calibration,
the native-to-ZF bridge, and the Foundation/Choice dependencies in the
retained predecessor. Relevant inherited theorem interfaces include
\path{LogicalSeed246.halts_iff_native_refutation},
\path{WFTheory.native_refutation_iff_zf}, and
\path{ConstructibleChoice.zf_refutation_iff_zfc_refutation}.
The ordinary interpretation arguments and their exact source excerpts
are included through the pinned 236-state dependency. The existing
external 240-state Lean endpoint is inherited; no new 235-state kernel
endpoint is asserted.

All new ingredients concern explicit finite objects and induction over
natural numbers or finite lists. They can be formalized in ZFC. Positive
durations ensure that an infinite source run gives infinitely many
machine transitions, and a finite halting source run gives a finite
halting machine run. Conversely, no runtime block has an interior halt
except the specified halt instruction. Negating the displayed chain
proves the main nonhalting biconditional.

No assertion about external standardness of an arbitrary model of ZFC
is needed. The proof concerns the arithmetized machine and refutation
predicates interpreted inside ZFC; the ordinary metatheorem about the
standard blank run is then obtained in the usual ambient mathematics.
If ZFC is consistent, $M_{235}$ actually does not halt, and G\"odel's
second incompleteness theorem prevents ZFC from proving that
nonhalting statement. Mere consistency is not used to claim the stronger
two-sided independence conclusion: unprovability of a false halting
assertion requires a suitable additional soundness hypothesis.
